\chapter{自动化估值建模}
\chapterlead{估值的结论是观点，但从假设到结论的每一步都应该是能重算的。}

\section{估值链}

卖方研报的盈利预测通常从业务驱动出发：白酒拆成茅台酒与系列酒的“吨数 × 吨价”，动力电池拆成“GWh × 单价”。
如果估值模型直接写一个合并收入增速，研报里讲的量价故事和估值里的数字之间就没有可以检查的对应关系。
\sysname{} 的估值链因此按“分部预测 → 三表联动 → 现金流折现 → 股权价值桥 → 程序评级”的顺序搭建（图~\ref{fig:val-chain}），
每一步的输入都来自上一步的输出或假设文件，不允许跳步。

\begin{figure}[htbp]
\centering
\resizebox{\linewidth}{!}{\begin{tikzpicture}[node distance=4mm and 5mm, every node/.style={font=\small}]
  \node[prog,fill=progc!18,minimum width=1.7cm] (tt) {真值表};
  \node[doc,below=7mm of tt,text width=2.6cm,minimum height=9mm] (as) {估值假设\\ \tiny valuation\_assumptions.json};
  \node[prog,right=7mm of tt,minimum width=2.3cm] (seg) {分部量价预测\\ \scriptsize 销量 × 单价、毛利率};
  \node[prog,right=of seg,minimum width=2.2cm] (ts) {三表联动\\ \scriptsize 三个情景各一遍};
  \node[prog,right=of ts,minimum width=2.1cm] (dcf) {FCFF DCF\\ \scriptsize 三情景 + 敏感性};
  \node[prog,right=of dcf,minimum width=1.9cm] (bridge) {股权价值桥\\ \scriptsize 每股价值};
  \node[gate,below=7mm of bridge] (rate) {评级};
  \node[prog,below=7mm of dcf,minimum width=2.1cm] (rdcf) {反向 DCF};
  \node[prog,below=7mm of ts,minimum width=2.2cm] (comps) {可比公司};
  \node[doc,below=7mm of seg,text width=2.0cm,minimum height=9mm] (xlsx) {公式联动工作簿\\ \tiny model.xlsx};
  \draw[arr] (tt) -- (seg); \draw[arr] (as) -| ([xshift=-3mm]seg.south); \draw[arr] (seg) -- (ts); \draw[arr] (ts) -- (dcf); \draw[arr] (dcf) -- (bridge);
  \draw[arr] (bridge) -- (rate); \draw[arr] (dcf) -- (rdcf); \draw[arr] (ts) -- (comps);
  \draw[darr] (comps) -- (xlsx); \draw[darr] (rdcf.south) |- ([yshift=-5mm]xlsx.south) -- (xlsx.south);
\end{tikzpicture}}
\caption{估值链。DCF 的每个情景向量都逐期取自三表模型的同一情景。}
\label{fig:val-chain}
\end{figure}

\subsection{分部量价预测}

分部的基期收入、成本和销量取自经营数据表（已经逐格回查原文），程序算出基期单价和毛利率。分析师在假设文件里为每个分部、每个情景给出
销量增速、单价增速和毛利率变动，程序递推出各期分部收入，加总得到合并收入和综合毛利率。没有拆分的部分记为“其他业务”，预测期保持不变。
费用率拆成税金及附加、销售、管理、研发和其他五项，可以按情景分别覆盖。

这样一来，研报里写“系列酒量价承压”，就必须在系列酒那一行的销量或单价增速上体现出来；主笔要写收入增速多少，只能引用程序从量价推出的结果。
销量单位因行业而异，宁德时代写 GWh、茅台写吨，单位只影响工作簿的表头，不进入计算。

\subsection{三表联动}

三表引擎在三个情景下各跑一遍。资本开支按收入比例，折旧作用于期初固定资产加半年资本开支，应收、存货、应付按周转天数，利息按平均有息负债，
现金是资产负债表的平衡项。默认驱动全部取基期实际值，分析师只需要覆盖他有依据的那几个。每个预测期都检查资产负债表是否平衡。

这里有一个口径上的取舍。早期版本让分析师在假设文件里直接写 DCF 用的 EBIT 利润率，三表模型另跑一套，结果研报首页的盈利预测表（来自三表）
和估值（来自 DCF）用的是两套利润率，复核编辑无法判断哪个对。现在 DCF 的情景向量全部取自三表模型的同一情景，
程序逐期核对两边的收入、EBIT、折旧摊销、资本开支、营运资本，偏差超过 $10^{-6}$ 百万元就把模型状态置为“阻断”。
三家真实公司运行时的最大偏差在 $10^{-11}$ 量级。要让悲观情景利润率更低，分析师必须去改毛利率或费用率，利润率的变化因此总能落到一个可以解释的经营驱动上。

\subsection{现金流折现与股权价值桥}

DCF 采用企业自由现金流（FCFF），预测期现金流按年中折现，终值同时用永续增长法和退出倍数法计算，并互相报告隐含参数：
永续增长法隐含的退出倍数，以及退出倍数法隐含的永续增长率。折现率为加权平均资本成本（WACC），股权成本按 CAPM 加规模、公司特定与国别溢价，
每个情景可以单独调整。计算前程序会拒绝不合法的输入，例如任一情景的 WACC 不大于永续增长率（此时 Gordon 公式不收敛）。

企业价值到每股价值经过一座逐项列示的股权价值桥：加现金、非经营性资产、联营投资，减有息负债、租赁负债、少数股东权益等，再除以摊薄股本。
EBIT 采用经营性口径（利润总额加财务费用，减投资收益和公允价值变动收益），被剔除的投资类资产需要在假设文件里作为非经营性资产加回，
否则报告会提示估值偏保守。

\subsection{敏感性、情景与检查}

基准情景上生成三张 $5\times5$ 的敏感性矩阵：WACC × 永续增长率、WACC × 退出倍数、收入增速 × EBIT 利润率。程序检查方向是否合理：
每股价值应随 WACC 递减、随永续增长率和利润率递增。我们有意不要求“每股价值随收入增速递增”，因为增长需要投入营运资本和资本开支，
当投入资本回报率低于 WACC 时，增长反而毁灭价值。样例公司的矩阵里就出现了这种情况：利润率下调 3 个百分点的那一行，收入增速越高，每股价值越低。

此外程序还检查情景排序（悲观 $\le$ 基准 $\le$ 乐观）、终值占企业价值是否超过 75\%、预测期末利润率较历史提升是否超过 3 个百分点等，
不满足时给出警示，严重的直接阻断。

\section{程序评级}

评级按基准每股价值相对参考股价的上行空间 $u$ 由程序给出：$u\ge20\%$ 买入，$u\ge5\%$ 增持，$u\ge-5\%$ 中性，$u\ge-20\%$ 减持，其余卖出。
阈值可以在假设文件里改，改后的规则文字随评级一起输出。

把评级交给程序，是因为评级必须和估值数字一致，而程序评级让这种一致性按构造成立，护栏只需核对正文里的评级字样与产物是否相同。
代价是评级只是基准情景的函数，没有考虑情景偏度和可比公司。我们的处理是把这些交叉验证并列放在研报的估值章节，由主笔解释评级与它们的关系，
但主笔不能改评级本身。如果分析师认为评级不合理，唯一的办法是修改并书面说明假设，然后重跑。

\section{反向 DCF：现价隐含了什么}

正向 DCF 回答“按我们的假设值多少钱”，反向 DCF 回答“要支撑现在的股价，需要什么样的假设”。程序对四个杠杆分别求解：
收入复合增速、平均 EBIT 利润率、隐含回报率（使每股价值等于股价的折现率）和永续增长率，每次只动一个，其他保持基准，并把解与公司历史最好水平比较。
表~\ref{tab:val-rdcf} 是仓库内三家公司的结果。

\begin{table}[htbp]
\centering
\caption{反向 DCF（基准假设 → 现价隐含值）。数据为仓库内回归运行所用的假设文件，与交付研报所用的工作区假设不同。}
\label{tab:val-rdcf}
\footnotesize
\begin{tabularx}{\linewidth}{@{}l c L L L L@{}}
\toprule
\hd{公司} & \hd{股价 / 基准价值} & \hd{收入复合增速} & \hd{平均 EBIT 利润率} & \hd{隐含回报率} & \hd{永续增长率} \\
\midrule
贵州茅台 & 1237.0 / 1531.4 & 4.32\% → $-1.91\%$ & 67.34\% → 53.78\% & 7.65\% → 9.03\% & 2.50\% → 0.73\% \\
\rowband
宁德时代 & 291.11 / 383.31 & 12.75\% → 4.09\% & 16.44\% → 12.27\% & 8.31\% → 10.64\% & 2.50\% → $-0.67\%$ \\
招商轮船 & 20.15 / 7.51 & 4.02\% → 25.91\%（超出历史最好 9.22\%） & 28.64\% → 52.48\%（超出历史最好 25.12\%） & 8.36\% → 4.53\% & 1.50\% → 5.75\% \\
\bottomrule
\end{tabularx}
\end{table}

茅台和宁德时代的现价低于基准价值，四个杠杆的隐含值都不高于基准，说明现价计入的预期比我们的基准更保守。招商轮船正好相反：
现价要求预测期收入年均增长 25.9\% 或平均利润率 52.5\%，都远超公司历史最好水平。这个结果在模型内部是自洽的，但它同时暴露了 DCF 用于强周期公司的局限：
以 2025 年为基期外推的现金流，几乎把当前运价周期的景气全部折掉了，周期股的定价更多反映周期高点的盈利与重置成本。
程序不会替我们修正这一点，我们要求主笔把“估值方法对周期股的适用性”写进研报的风险与局限一节。

\section{可比公司估值}

可比公司模块计算 EV/营业收入、EV/EBITDA、EV/EBIT、P/E 等倍数，分母不为正的记为“无意义”不参与统计。同业分为核心、次要和剔除三层，
剔除必须写理由，只有核心同业驱动估值区间。以中位数为中枢、四分位为区间，核心样本不少于 4 家时用四分位距标记离群值，
标记后在统计中剔除，但数据保留在表里，读者可以看到被剔除的是谁。

倍数的溢价或折价常常由利润率差异解释，所以程序同时计算目标公司与核心同业的 EBITDA、EBIT 和净利率差。
目标公司的 EV 用与 DCF 相同的股权价值桥，保证两种方法口径一致。程序最后报告 DCF 基准相对可比公司中枢的偏离，以及两个区间是否重叠。

\section{公式联动工作簿}

评审和研究员更习惯在 Excel 里检查估值，所以工作簿必须是“活的”。\file{model.xlsx} 共 15 个工作表（表~\ref{tab:val-sheets}），由标准库手写 OOXML 生成，
不依赖任何第三方 Excel 库。DCF 页的税后经营利润、FCFF、折现因子、终值、股权价值桥和每股价值都是公式；三表页也写成公式，
资产负债表平衡检查本身也是公式。“三表-假设”页有一个情景开关，改成 1、2、3 即可在悲观、基准、乐观之间切换，三表与 DCF 联动重算，
结果与引擎各情景的输出一致。输入单元格为蓝字，公式为黑字，跨表引用为绿字，沿用卖方模型的惯例。

\begin{table}[htbp]
\centering
\caption{估值工作簿的页面结构（以中际旭创交付工作簿为例）}
\label{tab:val-sheets}
\small
\begin{tabularx}{\linewidth}{@{}l L l L@{}}
\toprule
\hd{工作表} & \hd{内容} & \hd{工作表} & \hd{内容} \\
\midrule
估值汇总 & 三情景每股价值、评级、关键假设 & 三表-假设 & 情景开关与各情景驱动 \\
\rowband
盈利预测 & 研报首页的预测表 & 三表-利润表 & 公式联动 \\
DCF & FCFF、折现、终值、价值桥 & 三表-资产负债表 & 公式联动、平衡检查 \\
\rowband
WACC & 股权成本、债务成本、权重 & 三表-现金流量表 & 公式联动 \\
敏感性 & 三张 $5\times5$ 矩阵 & 历史财务 & 真值表历史期 \\
\rowband
非经营性资产 & 加回明细 & 分季度 & 单季拆分 \\
分部预测 & 分部量价递推 & 经营数据 & 经营数据表 \\
\rowband
模型检查 & 勾稽、同口径、方向性检查结果 & & \\
\bottomrule
\end{tabularx}
\end{table}

质检侧另有一个只读的工作簿审计工具，检查三类问题：结构（公式数、缺缓存值、错误值、失效引用）、公式控制（跨期公式“孤岛”，即中间某一年被手工改成数值；
公式里写死的数字；易变函数）、五组勾稽（资产负债平衡、现金衔接、留存收益、固定资产、有息负债的滚动）。在 86 本历史工作簿上没有发现公式孤岛。

\section{自研引擎与参照实现的一致性}

估值引擎参照了 OpenAI Public Equity Investing 插件（0.1.31）的计算口径。该插件的许可证标注为仅限内部使用，因此我们没有引入它的任何代码，引擎全部自研，参照关系在 \file{THIRD\_PARTY.md} 中逐项列明。
为了证明“口径一致”，我们设了两层守护。第一层是对照测试：本地存在原插件时，在随机生成的输入上逐项比较，三表引擎 64 组、情景表 180 组以上、
DCF 21 组 × 3 情景、业绩分析 300 组以上，三表与情景表要求浮点逐位相等，DCF 要求相对误差小于 $10^{-9}$。
第二层是快照测试：不需要原插件，用样例跑完整流水线后与 \file{tests/golden/} 中的四个快照逐字段比较，共 1,712 个叶子字段全部相等。

逐位相等的要求让写代码更费事，例如 \code{(begin+0.5*capex)*rate} 不能写成 \code{begin*rate+0.5*capex*rate}，因为浮点运算顺序不同结果就差最后一位。
好处是任何差异都是真差异，不必讨论“差 $10^{-7}$ 算不算一致”。2026 年 9 月 28 日我们有意修改了几处口径（经营性 EBIT、三表驱动 DCF、可比公司增加利润率），
重新生成了快照，并在快照说明中记录了差异原因。

\begin{note}[已知的口径简化]
折现起点是最近财年末而不是估值基准日，终值按第 $N$ 年末折现；流动性压力测试沿用原插件口径 FCF ×（1+冲击），FCF 为负时方向相反；
三表模型期初留存收益置 0、权益合并列示，只影响列示，不影响勾稽。估值链只覆盖非金融企业，银行和保险的负债是经营原料，FCFF 无定义。
这些简化都写在估值报告的口径说明里。
\end{note}
